% ==============================================================================
% METODOLOGÍA
% ==============================================================================
% Instrucciones de la plantilla:
% - Presentar de manera clara y estructurada el proceso seguido durante la práctica de laboratorio.
% - Comenzar con la configuración del entorno de trabajo (herramientas, lenguajes, bibliotecas: OpenGL, GLU, Assimp, GLEW, GLFW, GLM, etc.).
% - Detallar los algoritmos gráficos implementados explicando su fundamento teórico y su adaptación práctica (dibujo de formas, transformaciones geométricas, etc.).
% - Incluir imágenes con títulos y descripciones de los pasos ejecutados.
% - Suficientemente detallado para que cualquier lector pueda replicar el programa.
% - REGLA: "Una metodología por actividad".
% ==============================================================================

\section{Metodología}

\subsection{Configuración del entorno de trabajo}
Para el desarrollo de la práctica se utilizó un entorno de desarrollo estructurado en C++ configurado con soporte para OpenGL moderno. En la Tabla \ref{tab:requerimientos} se detallan las herramientas, bibliotecas y componentes del sistema empleados.

\begin{table}[H]
\centering
\renewcommand{\arraystretch}{1.25}
\begin{tabular}{|p{4.2cm}|p{10.5cm}|}
\hline
\textbf{Elemento} & \textbf{Especificación / Versión empleada} \\ \hline
Sistema operativo & Linux / Windows 11 (64 bits) \\ \hline
Entorno de desarrollo & Visual Studio Code / Microsoft Visual Studio \\ \hline
Compilador y estándar & GCC / MSVC (soporte para \texttt{C++20}) \\ \hline
API gráfica & OpenGL 3.3 (Perfil Core) \\ \hline
Bibliotecas principales & GLFW (manejo de ventanas y eventos), GLAD (cargador de funciones), GLM (matemáticas vectoriales y matriciales), Assimp (carga de modelos 3D) \\ \hline
Hardware de prueba & Tarjeta gráfica compatible con OpenGL 3.3 Core \\ \hline
\end{tabular}
\caption{Requerimientos de hardware, software y bibliotecas utilizadas.}
\label{tab:requerimientos}
\end{table}

\subsection{Actividad 1: Construcción e Integración del Escenario Completo 3D}

En esta actividad se diseñó un entorno tridimensional mediante el modelado de un edificio y elementos complementarios en Blender, para posteriormente exportarlo al formato FBX e integrarlo en una aplicación gráfica desarrollada en C++ utilizando la API de OpenGL y la biblioteca Assimp.

\subsubsection{Algoritmo y fundamento matemático}
Explicar el fundamento teórico del algoritmo (por ejemplo, el algoritmo de rasterización de líneas de Bresenham o cálculo de transformaciones afines):
\begin{equation}
    P' = M_{proyecci\acute{o}n} \cdot M_{vista} \cdot M_{modelo} \cdot P
    \label{eq:transformacion}
\end{equation}

\subsubsection{Implementación del código}
A continuación se muestra el segmento clave de código que implementa la solución:

\begin{lstlisting}[language=C++, caption={Configuración de buffers (VAO, VBO) y dibujo de geometrías.}, label={lst:actividad1}]
// Generación y enlace del Vertex Array Object (VAO) y Vertex Buffer Object (VBO)
unsigned int VAO, VBO;
glGenVertexArrays(1, &VAO);
glGenBuffers(1, &VBO);

glBindVertexArray(VAO);
glBindBuffer(GL_ARRAY_BUFFER, VBO);
glBufferData(GL_ARRAY_BUFFER, sizeof(vertices), vertices, GL_STATIC_DRAW);

// Definición de atributos de vértices (posición xyz)
glVertexAttribPointer(0, 3, GL_FLOAT, GL_FALSE, 3 * sizeof(float), (void*)0);
glEnableVertexAttribArray(0);
\end{lstlisting}

\subsection{Actividad 2: Variación del Escenario 3D}
% Describir aquí la metodología correspondiente a la Actividad 2.
% Recordar: "Una metodología por actividad".
Esta actividad consistió en implementar una variación al escenario original agregando más elementos geométricos en sus alrededores desde Blender, modificando la estructura del modelo e integrando un personaje navegable.